\section{Asignación de cómputo: el modelo de 1T como boleto de entrada y la proporción 3:1:1}

Uno de los juicios más concretos sobre recursos en la entrevista es la idea de que «un modelo de 1T es el boleto de entrada», junto con el cambio en la proporción de cómputo. Luo Fuli (罗福莉) considera que, si se quiere alcanzar un nivel de Agent cercano al de Claude Opus 4.6, un modelo base de escala 1T puede ser el boleto de entrada. Sin embargo, al entrar en la etapa de Agent, la asignación de cómputo ya no puede quedar acaparada por el preentrenamiento.

\begin{figure}[H]
\centering
\includegraphics[width=0.92\textwidth]{compute-allocation.png}
\caption{Tendencia de la asignación de cómputo desde la etapa de Chat hasta la etapa de Agent.}
\end{figure}

\begin{knowledgebox}{Lectura de la figura: qué significan 3:5:1 y 3:1:1}
Las proporciones 3:5:1 y 3:1:1 de la figura son cifras mencionadas de forma oral en la entrevista y sirven para expresar el desplazamiento del peso de los recursos. En la era de Chat, el preentrenamiento tenía una participación mayor; en la era de Agent, la relación entre research, pre-train y post-train se acerca más al equilibrio, y la proporción de GPU destinada al posentrenamiento aumenta de forma notable.
\end{knowledgebox}

\subsection{En qué sentido 1T es un boleto de entrada}

«Un modelo de 1T es el boleto de entrada» no significa que basten 1T de parámetros, ni que los modelos pequeños carezcan de valor. Apunta a un umbral de capacidad para un Agent de frontera: si se pretende competir con los sistemas de Agent más fuertes, el modelo base debe alcanzar un nivel suficientemente alto en escala, conocimiento, razonamiento, código y capacidad para tareas de varios pasos. El posentrenamiento y el framework pueden amplificar el modelo, pero no pueden crear de la nada todas las capacidades.

\begin{warningbox}{La escala no es la única respuesta}
El número de parámetros es solo una variable de entrada. Un sistema de Agent también depende de los datos de entrenamiento, el contexto, la capacidad de programación, el uso de herramientas, el posentrenamiento, el entorno de la tarea, el costo y la ejecución organizacional. Interpretar «1T como boleto de entrada» como «basta con acumular parámetros para ir a la cabeza» es una lectura errónea.
\end{warningbox}

\subsection{Equilibrio entre post-train y pre-train}

Luo Fuli considera que, en los equipos de primer nivel, la proporción de recursos entre preentrenamiento y posentrenamiento ya se acerca a 1:1. Esto significa que el posentrenamiento ya no es el cierre que sigue al preentrenamiento, sino un frente principal al mismo nivel que este. En las tareas de Agent en particular, el posentrenamiento determina directamente si el modelo puede adaptarse al entorno, a las herramientas y a las tareas de largo alcance.

\begin{importantbox}{La asignación de recursos es una decisión estratégica}
Cómo repartir las GPU no es una cuestión financiera, sino una decisión sobre la ruta técnica. Seguir concentrando todas las GPU en el preentrenamiento, o reservar suficientes para research, post-train, rollout y eval, produce capacidades organizacionales completamente distintas.
\end{importantbox}

\subsection{Resumen de la sección}

En la era de Agent, la asignación de cómputo es más equilibrada y más compleja. La competencia de frontera requiere tanto un modelo base grande como recursos para el posentrenamiento y la investigación. El mensaje central de la proporción 3:1:1 es que el posentrenamiento pasó de ser un actor secundario a ser el protagonista.
